\documentclass[a4paper]{article}
% Español (México) con XeLaTeX: fontspec para fuentes Unicode, polyglossia
% para la división silábica y los nombres del documento en la variante mexicana.
\usepackage{fontspec}
\usepackage{polyglossia}
\setdefaultlanguage[variant=mexican]{spanish}
% Los nombres chinos van como «pinyin (original)»: xeCJK da a los caracteres
% Han su propia fuente; sin ella XeLaTeX los omite con un aviso.
% FandolSong no cubre algunos caracteres de nombres (U+73FA); WenQuanYi Zen Hei sí.
\usepackage[AutoFallBack=true]{xeCJK}
\setCJKmainfont{FandolSong-Regular.otf}
\setCJKfallbackfamilyfont{\CJKrmdefault}{WenQuanYi Zen Hei}
% polyglossia no traduce los nombres de listings.
\AtBeginDocument{\providecommand{\lstlistingname}{}\renewcommand{\lstlistingname}{Listado}%
  \providecommand{\lstlistlistingname}{}\renewcommand{\lstlistlistingname}{Índice de listados}}
\usepackage{amsmath, amssymb}
\usepackage{graphicx}
\usepackage[margin=2.35cm]{geometry}
\usepackage[most]{tcolorbox}
\usepackage{booktabs}
\usepackage{tabularx}
\usepackage{longtable}
\usepackage{array}
\usepackage{float}
\usepackage{xcolor}
\usepackage{hyperref}
\usepackage{fancyhdr}
\setCJKmainfont[AutoFakeBold=2.4]{AR PL UMing CN}
\setCJKsansfont[AutoFakeBold=2.4]{Droid Sans Fallback}
\setCJKmonofont[AutoFakeBold=2.4]{Droid Sans Fallback}
\hypersetup{colorlinks=true, linkcolor=blue!55!black, urlcolor=blue!60!black}
\graphicspath{{./}{figures/}}
\setlength{\parskip}{0.55em}
\setlength{\parindent}{2em}
\setlength{\emergencystretch}{3em}
\renewcommand{\arraystretch}{1.18}
\newtcolorbox{knowledgebox}[1]{enhanced, colback=blue!5!white, colframe=blue!70!black, colbacktitle=blue!70!black, coltitle=white, fonttitle=\bfseries, title={#1}, sharp corners, breakable}
\newtcolorbox{importantbox}[1]{enhanced, colback=yellow!10!white, colframe=yellow!75!black, colbacktitle=yellow!75!black, coltitle=black, fonttitle=\bfseries, title={#1}, sharp corners, breakable}
\newtcolorbox{warningbox}[1]{enhanced, colback=red!5!white, colframe=red!70!black, colbacktitle=red!70!black, coltitle=white, fonttitle=\bfseries, title={#1}, sharp corners, breakable}
\pagestyle{fancy}
\fancyhf{}
\fancyfoot[C]{\thepage}
\fancyfoot[R]{\footnotesize\color{gray}\href{https://github.com/hqhq1025/ai-course-notes}{AI Course Notes}}
\renewcommand{\headrulewidth}{0pt}
\renewcommand{\footrulewidth}{0pt}
\newcommand{\videourl}{https://www.youtube.com/watch?v=858HR43pegk}
\newcommand{\repourl}{https://github.com/hqhq1025/ai-course-notes}

\begin{document}
\begin{titlepage}
\centering
{\Large Notas didácticas de revisión técnica\par}
\vspace{1.0cm}
{\huge\bfseries Kimi Linear, Sparse Attention y la arquitectura de modelos de la próxima generación}\par
\vspace{0.5cm}
{\Large Zhang Xiaojun conversa con Yang Songlin (杨松琳)\par}
\vspace{0.35cm}
{\large Elaborado a partir del video público de Zhang Xiaojun Podcast, la transcripción generada en GPU, la estructura de capítulos y diagramas conceptuales propios\par}
\vspace{0.3cm}
{\large 2025-11-03\par}
\vspace{0.85cm}
\includegraphics[width=0.88\textwidth,height=0.42\textheight,keepaspectratio]{cover.jpg}\par
\vfill
\begin{tcolorbox}[width=0.92\textwidth, colback=black!2!white, colframe=black!55, sharp corners]
\textbf{Título del video}: 119. Kimi Linear, ¿Minimax M2? Yang Songlin repasa la historia de las variantes algorítmicas y anticipa propuestas de mejora para las arquitecturas futuras\par
\textbf{Canal del video}: Zhang Xiaojun Podcast\par
\textbf{Duración del video}: 01:43:26\par
\textbf{Enlace del video}: \href{\videourl}{\nolinkurl{\videourl}}\par
\textbf{Notas sobre el material}: YouTube no ofrece subtítulos disponibles; el texto se elaboró a partir de la transcripción con faster-whisper large-v3 en CUDA, los capítulos del video, la portada y figuras didácticas propias. Las tomas fijas de la entrevista no se reutilizan en el cuerpo del texto.\par
\textbf{Página del proyecto}: \href{\repourl}{\nolinkurl{\repourl}}\par
\end{tcolorbox}
\end{titlepage}

\tableofcontents
\newpage
\section{Introducción: por qué en 2025 se volvió a «tallar la Attention»}

Esta sección establece primero el contexto técnico de todo el episodio. Yang Songlin (杨松琳) habla de uno de los módulos más importantes de la arquitectura básica de los modelos grandes: la Attention. En los últimos años, la parte FFN del Transformer ya se modificó a fondo con MoE (Mixture of Experts, mezcla de expertos), y modelos como DeepSeek hicieron de MoE un consenso. La siguiente pieza que podría «tallarse» es la Attention. El motivo es directo: el contexto largo, el CoT largo, la Agentic AI y el decoding de textos largos pusieron en primer plano la carga de cómputo y de KV cache de la Full Attention.

Este episodio busca responder cuatro preguntas. Primera: ¿en qué consiste exactamente lo «linear» de la Linear Attention? Segunda: ¿por qué es importante el módulo KDA de Kimi Linear? Tercera: ¿qué compromisos implican, respectivamente, la atención lineal híbrida de Kimi, la atención dispersa de DeepSeek y el regreso de MiniMax M2 a la Full Attention? Cuarta: ¿por qué la investigación en arquitecturas no puede limitarse a la elegancia matemática y debe coordinarse con la afinidad con el hardware, los algoritmos paralelos, la multiplicación de matrices y la optimización de kernels?

\begin{importantbox}{Tesis central del episodio}
Cuando el crecimiento de los datos se desacelera, aumenta la demanda de contexto largo y el costo de inferencia se vuelve más pesado, la innovación en arquitecturas vuelve a ser decisiva. Las nuevas líneas de Attention no buscan sustituir el nombre Transformer, sino obtener, con los mismos FLOPs, una loss más baja, menos carga de KV cache y un decoding más eficiente en secuencias largas.
\end{importantbox}

\begin{knowledgebox}{Nota sobre la estrategia visual}
Este video es una entrevista con encuadre fijo, sin slides ni proyección de artículos. El texto no incluye fotogramas repetidos de los participantes; todas las figuras son diagramas conceptuales elaborados para esta nota, que explican la complejidad de la atención, KDA, las tres líneas de Attention, la Scaling Ladder, la arqueología de algoritmos y la afinidad con el hardware.
\end{knowledgebox}

\subsection{Resumen de la sección}

EP119 es una revisión panorámica de arquitecturas, no una entrevista común. Reúne en un mismo mapa las distintas apuestas que las empresas de modelos hicieron en 2025 sobre la Attention: Kimi apuesta por Linear/Hybrid, DeepSeek por Sparse y MiniMax M2 regresa por ahora a la Full Attention.
\section{Mapa de problemas de Attention: CoT largo, KV Cache y Decoding}

La sección anterior explicó por qué hay que volver a investigar Attention; esta sección empieza por desglosar el problema. La ventaja de Full Attention es su gran capacidad expresiva, su entrenamiento estable y la amplia experiencia acumulada con ella; su desventaja es que, con secuencias largas, tanto el cómputo como la memoria de GPU resultan costosos. El CoT largo y la Agentic AI generan cadenas de razonamiento de decenas de miles de tokens; el decoding genera token por token, el KV cache crece cada vez más y el costo de inferencia se convierte en un cuello de botella real.

\begin{figure}[H]
\centering
\includegraphics[width=0.96\textwidth]{attention-problem-map.png}
\caption{Mapa de problemas de Attention: la inferencia con contexto largo lleva el KV cache, el decoding y la atención global hasta el cuello de botella. Diagrama conceptual propio, elaborado a partir de la conversación de 00:12:20--00:14:39.}
\end{figure}

\begin{knowledgebox}{Lectura de la figura: un contexto largo no solo aumenta la longitud de la entrada}
El CoT largo hace que el modelo produzca tokens de forma continua durante la fase de generación. En cada paso, Full Attention debe interactuar con el contexto histórico, el KV cache debe guardar las key/value históricas y el costo del decoding aumenta a medida que el contexto se alarga. El objetivo de las nuevas variantes de Attention es reducir estos costos sin una pérdida notable de desempeño.
\end{knowledgebox}

\subsection{Complejidad aproximada de Full Attention}

En la Softmax Attention autorregresiva estándar, dada una longitud de secuencia \(L\), el modelo construye \(QK^\top\) para obtener una matriz de puntuaciones de atención de \(L \times L\), y después aplica la causal mask, el softmax y la multiplicación por \(V\). Por ello, el cálculo de atención en la fase de entrenamiento o de prefill suele tener complejidad cuadrática:
\[
\text{cost}_{\text{full}} = O(L^2).
\]
Aquí, \(L\) es la longitud de la secuencia; \(Q,K,V\) son, respectivamente, las representaciones query, key y value; \(O(L^2)\) indica que el costo crece con el cuadrado de la longitud de la secuencia.

\begin{warningbox}{La complejidad no es el único indicador}
Full Attention tiene una complejidad teórica alta, pero está compuesta por grandes multiplicaciones de matrices, por lo que se adapta muy bien al hardware. Muchos algoritmos lineales o dispersos tienen menor complejidad teórica, pero si su paralelismo es deficiente, sus kernels son difíciles de escribir o su acceso a memoria es poco eficiente, su velocidad real no necesariamente es mejor.
\end{warningbox}

\subsection{Por qué importa el KV Cache}

El KV cache es la caché que guarda las key/value históricas durante la inferencia. Permite que el modelo no tenga que recalcular todas las representaciones históricas al generar un nuevo token, pero a cambio la memoria de GPU crece con el número de capas, la longitud de la secuencia, el batch size y la dimensión oculta. Una ventaja de la ruta de atención lineal híbrida es que muchas capas pueden funcionar como el estado recurrente de una RNN, lo que reduce el KV cache; la ruta de atención dispersa, en cambio, reduce el cómputo de cada paso sobre todo al atender solo a los tokens Top-K o a los de una ventana.

\begin{knowledgebox}{Cuentas aproximadas del KV cache}
Si un modelo tiene \(N\) capas, cada capa guarda key y value, la longitud de la secuencia es \(L\), el batch size es \(B\) y la dimensión de la representación KV de cada token es \(d\), entonces el orden de magnitud del KV cache puede escribirse como:
\[
\text{KV cache} \propto 2 \times N \times B \times L \times d.
\]
Aquí, el 2 inicial proviene de las dos cachés, una para key y otra para value. Esta expresión muestra que, en cuanto crece la longitud de contexto \(L\) o el batch size \(B\), la presión sobre la memoria de GPU aumenta linealmente; si todas las capas conservan el KV cache de Full Attention, la inferencia con contexto largo resulta muy costosa.
\end{knowledgebox}

\begin{importantbox}{Por qué ahorrar KV cache aumenta el rendimiento de procesamiento}
Cuando el KV cache es más pequeño, en la misma memoria de GPU cabe un batch size mayor, y un batch size mayor a su vez aumenta el rendimiento de procesamiento del servidor. Por ello, la ruta hybrid linear de Kimi y similares no solo reduce el uso de memoria de GPU, sino que también influye en la economía del servicio de inferencia.
\end{importantbox}

\subsection{Resumen de la sección}

El cuello de botella de Attention no es un problema exclusivo de FLOPs, sino el efecto conjunto del cómputo, la memoria de GPU, el batch size, la latencia del decoding y la eficiencia del hardware con secuencias largas. Solo al entender estas restricciones se comprende por qué Kimi, DeepSeek y MiniMax eligieron rutas distintas.
\section{Linear Attention: reescribir la matriz cuadrática como una recurrencia lineal}

Esta sección explica la intuición básica de Linear Attention. Lo «linear» no significa que el modelo se vuelva una función lineal, sino que su complejidad crece de forma casi lineal con la longitud de la secuencia. Lo habitual es eliminar o reformular la Softmax y, mediante transformaciones algebraicas, escribir la atención como una recurrence similar a la de una RNN. En cada paso se mantiene un estado, de modo que procesar una secuencia de longitud \(L\) cuesta aproximadamente \(O(L)\).

\begin{figure}[H]
\centering
\includegraphics[width=0.94\textwidth]{linear-attention-intuition.png}
\caption{Intuición de Linear Attention: la matriz cuadrática se reescribe como una recurrence lineal. Diagrama conceptual propio, elaborado a partir de la conversación entre 00:07:04 y 00:11:19.}
\end{figure}

\begin{knowledgebox}{Lectura de la figura: Linear Attention no consiste en escribir una capa menos}
Softmax Attention construye de forma explícita una matriz de \(L\times L\). Linear Attention intenta comprimir la información histórica en un estado recurrente, de modo que cada paso de generación solo actualiza ese estado. Así se reducen la complejidad y la presión sobre la KV cache, pero hay que rediseñar la capacidad expresiva y la estabilidad del entrenamiento.
\end{knowledgebox}

\subsection{Su lugar dentro del Transformer}

El Transformer se compone principalmente de bloques de Attention y FFN apilados de forma repetida. En los últimos años, la FFN se transformó con MoE; Linear Attention, en cambio, interviene directamente en el módulo de Attention. Forma parte de la investigación sobre arquitecturas de preentrenamiento, a diferencia del trabajo sobre el optimizer, los datos de preentrenamiento o el post-training. Hoy es más común una arquitectura hybrid: algunas capas conservan Full Attention y otras se sustituyen por Linear Attention.

\begin{importantbox}{Motivación de la arquitectura hybrid}
Una Linear Attention pura puede quedarse corta en capacidad expresiva; una Full Attention pura resulta demasiado costosa. La arquitectura hybrid usa unas pocas capas globales como respaldo de la capacidad expresiva y muchas capas lineales para reducir el costo; hoy es el compromiso más práctico.
\end{importantbox}

\subsection{Resumen de la sección}

El núcleo de Linear Attention es reescribir la complejidad y la representación del estado. Puede reducir el costo de inferencia con contextos largos, pero tiene que resolver la capacidad expresiva, la estabilidad del entrenamiento, el paralelismo en hardware y la validación a gran escala.
\section{Kimi Linear y KDA: actualización de la memoria con granularidad más fina}

Esta sección entra en KDA (Kimi Delta Attention), el módulo central de Kimi Linear. En el programa se menciona que KDA parte de trabajos anteriores como Gated DeltaNet y sustituye la decay de granularidad gruesa por una decay de granularidad más fina: antes, las distintas dimensiones de un mismo attention head compartían una sola tasa de decaimiento; ahora cada dimensión puede tener su propia tasa de decaimiento. La intuición es que el hidden state, que es limitado, exprese mejor memorias de distintas escalas temporales.

\begin{figure}[H]
\centering
\includegraphics[width=0.96\textwidth]{kimi-linear-kda.png}
\caption{Kimi Linear / KDA: Kimi Delta Attention aprovecha mejor la memoria con una decay de granularidad más fina. Diagrama conceptual propio, elaborado a partir de la conversación de 00:14:39--00:18:56.}
\end{figure}

\begin{knowledgebox}{Lectura de la figura: KDA es una nueva combinación de ideas anteriores}
Gated DeltaNet aporta la base, la decay de granularidad gruesa asegura la eficiencia, la decay de granularidad fina aumenta la capacidad expresiva, KDA reorganiza la actualización de la memoria y, al final, todo se integra en una arquitectura hybrid. Muchas arquitecturas nuevas no se inventan de la nada: recombinan mecanismos anteriores en el contexto de hardware nuevo y escalas nuevas.
\end{knowledgebox}

\subsection{Scaling Ladder: superar la escala pequeña antes de ampliarla}

Esta subsección explica cómo Kimi reduce el riesgo de prueba y error en la arquitectura. La atención lineal no se lleva directamente a un modelo grande en cuanto se escribe la fórmula; tiene que validarse escala por escala.

Internamente, Kimi utiliza una Scaling Ladder: un módulo se compara primero con Full Attention a escala pequeña, y solo si su desempeño es suficientemente bueno pasa a una escala mayor para seguir haciendo scale. El proceso se parece a superar los niveles de un juego y evita gastar cómputo en prueba y error directamente a la escala máxima. El diseño de KDA también se convirtió poco a poco en candidato después de varias rondas de selección entre formas de combinación y de ampliación de escala.

\begin{figure}[H]
\centering
\includegraphics[width=0.96\textwidth]{scaling-ladder.png}
\caption{Scaling Ladder: tras superar la escala pequeña se pasa a una escala mayor, con comparación frente a Full Attention. Diagrama conceptual propio, elaborado a partir de la conversación de 00:18:56--00:20:20.}
\end{figure}

\begin{knowledgebox}{Lectura de la figura: la arquitectura no se decide de improviso}
La selección a escala pequeña reduce costos; cuando las métricas alcanzan el nivel requerido, se pasa a una escala mayor y se vuelve a comparar con la línea base de Full Attention. Solo si el costo, el desempeño y el comportamiento de scaling resultan satisfactorios puede volverse candidato para producción.
\end{knowledgebox}

\subsection{La proporción de mezcla tres a uno}

A continuación se trata el problema de la proporción en la arquitectura hybrid. Decidir cuántas capas globales conservar y cuántas sustituir por capas lineales consiste, en esencia, en encontrar entre la capacidad expresiva y la eficiencia un punto de ingeniería que pueda escalarse.

En la entrevista se menciona que Kimi Linear inserta una capa de atención completa por cada tres capas de KDA, y que la proporción tres a uno se ha ido volviendo un consenso. En esencia, esta proporción usa Full Attention para mantener la capacidad expresiva global y KDA para reducir el costo de inferencia en la mayoría de las capas. No es un teorema matemático, sino un compromiso empírico surgido de los experimentos de ingeniería.

\begin{knowledgebox}{Asimilación de los mecanismos de KDA}
\begin{tabularx}{\textwidth}{p{0.22\textwidth}p{0.34\textwidth}X}
\toprule
Mecanismo & Función & Por qué importa \\
\midrule
Delta Rule & Actualiza el estado de memoria de forma recursiva & Da a la attention una capacidad de actualización de estado similar a la de una RNN. \\
Gating & Controla qué información se conserva y cuál se olvida & Permite que el modelo elija qué información histórica debe permanecer. \\
Decay & Tasa de decaimiento de la memoria & Distintas informaciones pueden tener distintas escalas temporales. \\
Fine-grained decay & Cada dimensión tiene su propio decaimiento & Aprovecha mejor un hidden state limitado. \\
Capas globales hybrid & Insertan Full Attention de forma periódica & Garantizan como respaldo la expresión global y la interacción de largo alcance. \\
\bottomrule
\end{tabularx}
\end{knowledgebox}

\begin{warningbox}{KDA no se reduce a «quitar la Softmax»}
Si solo se quita la Softmax, el modelo puede quedarse sin suficiente capacidad expresiva. El valor de KDA está en rediseñar la actualización del estado, el gating y el decaimiento para que las capas lineales conserven en lo posible la información histórica útil; las capas de Full Attention se encargan de completar la interacción global.
\end{warningbox}

\subsection{Resumen de la sección}

Lo esencial de Kimi Linear no es «sustituir todo por atención lineal», sino usar KDA dentro de una arquitectura hybrid para actualizar la memoria de forma eficiente y garantizar un nivel mínimo de desempeño con unas pocas capas globales. La Scaling Ladder es el mecanismo que permite llevar este diseño, paso a paso, hasta el entrenamiento a gran escala.
\section{Tres rutas de Attention: Kimi, DeepSeek y MiniMax}

Antes se habló de Kimi; esta sección compara las tres rutas. Kimi sigue la ruta Linear/Hybrid, DeepSeek sigue la de Sparse Attention y MiniMax M2 regresa de la Linear/Hybrid de M1 a Full Attention. Las tres decisiones responden a la misma pregunta: cómo reducir el costo del decoding con contextos largos sin sacrificar demasiada capacidad expresiva ni capacidad de generalización.

\begin{figure}[H]
\centering
\includegraphics[width=0.96\textwidth]{linear-vs-sparse-vs-full.png}
\caption{Tres rutas de Attention: las distintas apuestas de Kimi, DeepSeek y MiniMax. Diagrama conceptual de elaboración propia, basado en la conversación de 00:20:20--00:27:00.}
\end{figure}

\begin{knowledgebox}{Lectura de la figura: tres rutas para resolver el mismo problema}
Linear/Hybrid ahorra KV cache, Sparse reduce los token a los que se atiende en cada paso y Full Attention tiene gran capacidad expresiva, pero es costosa. Kimi eligió la atención lineal híbrida, DeepSeek se inclina por la dispersión dinámica y MiniMax M2 regresó a la atención global, lo que muestra que cada equipo pondera de manera distinta el costo, la estabilidad y el desempeño en las tareas.
\end{knowledgebox}

\subsection{Kimi frente a DeepSeek}

La ruta hybrid de Kimi conserva algunas capas de atención global y sustituye la mayoría de las capas por atención lineal; por ello reduce de forma notable el KV cache y, al mismo tiempo, mejora la eficiencia del decoding. DeepSeek Sparse Attention, en cambio, usa un indexer para seleccionar los Top-K token, con lo que disminuyen los token que participan en el cálculo de cada paso, aunque no necesariamente se reduce el almacenamiento de KV cache de cada capa. Ambas apuntan a la eficiencia con contextos largos, pero optimizan aspectos distintos.

\begin{warningbox}{No hay que tratar a Linear ni a Sparse como la opción más avanzada en términos absolutos}
Linear/Hybrid, Sparse y Full Attention tienen cada una sus escenarios. La capacidad expresiva, la recuperación de largo alcance, el KV cache, el batch size, la implementación en hardware y la estabilidad del entrenamiento influyen en la elección final. Que una empresa de modelos regrese a Full Attention no significa necesariamente que Linear haya fracasado; también puede indicar que las tareas actuales y la estabilidad requieren más expresividad global.
\end{warningbox}

\subsection{Por qué MiniMax M2 regresó a Full Attention}

En el programa se menciona que MiniMax M1 fue uno de los pioneros de la atención lineal híbrida a gran escala, pero M2 regresó a Full Attention. Una posible razón es que Linear Attention no es lo bastante estable en algunas tareas de multi-hop reasoning o de recuperación compleja; otra es que Full Attention sigue siendo más madura en certidumbre de ingeniería, en piso de expresividad y en experiencia de entrenamiento. La elección de arquitectura no depende solo de la velocidad de inferencia, sino del desempeño en conjunto.

\begin{knowledgebox}{Tabla comparativa de las tres rutas}
\begin{tabularx}{\textwidth}{p{0.22\textwidth}p{0.25\textwidth}p{0.25\textwidth}X}
\toprule
Ruta & Ventajas & Costos & Preguntas que conviene observar \\
\midrule
Full Attention & Gran capacidad expresiva, experiencia madura & Costo alto de KV cache y de contextos largos & Si es más estable en razonamiento de múltiples saltos y recuperación compleja. \\
Linear/Hybrid & Ahorra KV cache, favorable para el decoding & Debe compensar la capacidad expresiva y la interacción de largo alcance & CoT largos, rendimiento de razonamiento de Agent. \\
Sparse Attention & Menos token por paso, permite codiseño con el hardware & El indexer que selecciona los token es decisivo & Si logra seleccionar con precisión el contexto útil. \\
\bottomrule
\end{tabularx}
\end{knowledgebox}

\begin{importantbox}{La elección de arquitectura es una optimización multiobjetivo}
Las empresas de modelos no eligen «el algoritmo óptimo», sino que hacen una optimización multiobjetivo entre desempeño, costo, estabilidad, implementación en hardware, experiencia del equipo y escenarios de producto. Que distintas empresas regresen a rutas distintas es algo normal.
\end{importantbox}

\subsection{Resumen de la sección}

La divergencia entre las tres rutas muestra que Attention todavía no ha convergido. Kimi prefiere usar hybrid para ahorrar costos, DeepSeek busca la dispersión dinámica en codiseño con el hardware y MiniMax eligió regresar a la Full Attention, más estable. En el futuro podrían surgir soluciones combinadas.
\section{De MoE a Attention: por qué la innovación en arquitectura vuelve a ser importante}

Esta sección sitúa a Attention dentro de una historia más amplia de las arquitecturas. El Transformer contiene Attention, FFN y codificación posicional. La FFN ya fue transformada por MoE; la codificación posicional pasó de posiciones absolutas a posiciones relativas, RoPE, etc.; Attention, por su parte, sigue evolucionando por vías como Full, Linear, Sparse y Sliding Window. La tesis central de la entrevista es que MoE ha sido el mayor avance de arquitectura de los últimos años y que la siguiente pieza podría ser Attention.

\begin{figure}[H]
\centering
\includegraphics[width=0.96\textwidth]{moe-to-attention.png}
\caption{De MoE a Attention: la FFN ya fue esculpida en forma de MoE, y Attention se convierte en la siguiente región por esculpir. Diagrama conceptual de elaboración propia, a partir de la conversación de 00:58:08--01:02:48.}
\end{figure}

\begin{knowledgebox}{Lectura de la figura: el Transformer no es un bloque monolítico}
El Transformer es un apilamiento de Attention + FFN. Al convertirse la FFN en MoE, cambió la forma de usar el cómputo; Attention sigue siendo el cuello de botella del contexto largo, por lo que vías como Linear/Sparse intentan seguir «esculpiendo» la arquitectura.
\end{knowledgebox}

\subsection{FLOPs, loss y arquitectura}

Esta subsección convierte la elección de arquitectura en una función objetivo comparable. Las empresas de modelos no modifican la arquitectura por novedad, sino para obtener mejores resultados de entrenamiento e inferencia con el mismo presupuesto de cómputo.

El objetivo elemental detrás de la innovación en arquitectura es: con los mismos FLOPs, obtener una loss más baja. FLOPs significa floating point operations, es decir, la cantidad de operaciones de punto flotante. Una buena arquitectura no se limita a reducir el cómputo, sino que, con el mismo presupuesto, expresa funciones con mayor eficacia, aprovecha mejor los datos y se adapta con más eficiencia al hardware.

\begin{importantbox}{La función objetivo de la innovación en arquitectura}
Bajo las mismas restricciones de datos y cómputo, la innovación en arquitectura reduce el costo de entrenamiento/inferencia, aumenta la capacidad expresiva o mejora la eficiencia en hardware; lo ideal es mejorar las tres a la vez.
\end{importantbox}

\subsection{Por qué el agotamiento de los datos hace más importante la innovación algorítmica}

A continuación, se devuelven los algoritmos al contexto de las restricciones de la industria. Mientras los datos y el cómputo puedan seguir creciendo sin mayor reflexión, la importancia de la innovación en arquitectura queda fácilmente oculta; cuando ambos se encarecen, la eficiencia algorítmica vuelve al centro.

Al inicio del programa se menciona que los datos, los algoritmos y el cómputo son los tres motores que impulsan la inteligencia artificial. Cuando el crecimiento de los datos se desacelera y el cómputo se vuelve limitado o costoso, el valor de los algoritmos y de la arquitectura vuelve a aumentar. Las empresas chinas, con un cómputo más limitado que el de Estados Unidos, tienen precisamente por ello más incentivos para innovar en direcciones como MoE, Sparse Attention, Linear Attention y NoPE/RoPE, buscando «ahorrar cómputo sin perder rendimiento».

\begin{knowledgebox}{Cambio en las restricciones de los tres motores}
\begin{tabularx}{\textwidth}{p{0.18\textwidth}p{0.38\textwidth}X}
\toprule
Variable & Práctica principal en el pasado & Presión actual \\
\midrule
Datos & Seguir ampliando los datos de internet/código/multimodales & El crecimiento de los datos públicos de alta calidad se desacelera. \\
Cómputo & Acumular GPU y ampliar la escala de entrenamiento & Se agravan las restricciones de costo, suministro y energía. \\
Algoritmos & Mantener la base del Transformer con ajustes menores & Se necesita mayor eficacia con los mismos FLOPs. \\
\bottomrule
\end{tabularx}
\end{knowledgebox}

\subsection{Resumen de la sección}

Attention vuelve a ser el foco porque el contexto largo y el costo de inferencia llevan a la atención global a una posición de cuello de botella. MoE transformó la FFN; la siguiente etapa podría ser que Linear/Sparse/Hybrid Attention transforme la capa de atención.
\section{Arqueología de algoritmos: cómo los artículos antiguos vuelven a brillar con el hardware nuevo}

Yang Songlin (杨松琳) enfatiza la «arqueología»: leer artículos antiguos, entender mecanismos antiguos y después recombinarlos frente a problemas nuevos y hardware nuevo. Muchos mecanismos ya habían aparecido en 2016, 2020 y 2021, pero en su momento faltaban el hardware adecuado, algoritmos paralelos, implementaciones de código abierto o validación a gran escala, así que quedaron sepultados en el mar de la literatura. DataNet, Delta Rule, Gated Linear Attention y otros pertenecen a este tipo de pistas reactivadas.

\begin{figure}[H]
\centering
\includegraphics[width=0.96\textwidth]{algorithm-archaeology.png}
\caption{Ruta de la arqueología de algoritmos: artículos antiguos, mecanismos antiguos y hardware nuevo se recombinan. Diagrama conceptual propio, elaborado a partir de la conversación de 01:23:12--01:29:50.}
\end{figure}

\begin{knowledgebox}{Lectura de la figura: la arqueología no es nostalgia, sino la búsqueda de herramientas reutilizables}
Entre Transformer, Sparse, DeltaNet, Gated Linear y KDA no hay un progreso en línea recta, sino mecanismos antiguos que se reevalúan en un contexto nuevo. Muchas ideas no aparecen por primera vez, sino que por primera vez reúnen las condiciones para hacer scale up.
\end{knowledgebox}

\subsection{Por qué los trabajos antiguos quedan sepultados}

Esta sección explica por qué la «arqueología» no es un adorno. Muchos artículos antiguos no estaban equivocados; en su momento carecían de las condiciones de ingeniería que les permitieran convertirse en la corriente principal.

Hay tres posibles razones por las que un trabajo antiguo queda sepultado. Primera: los algoritmos complementarios y los kernels no están maduros, de modo que a otros les resulta difícil reproducirlo o escalarlo. Segunda: el hardware y las tareas de la época no lo necesitaban, así que su valor no se hizo visible. Tercera: el código abierto es difícil de usar y la comunidad no puede darle continuidad. Yang Songlin enfatiza hacer bien el código para que la técnica pueda transmitirse; en realidad, esto es un componente importante del impacto de la investigación en arquitecturas.

\begin{warningbox}{Que funcione a pequeña escala no significa que funcione a gran escala}
Existen muchas variantes de arquitectura, y muchos algoritmos funcionan en tareas a pequeña escala, pero fallan en el preentrenamiento a gran escala. El Scaling Ladder, la implementación en hardware y la reproducción en código abierto son mecanismos importantes para evitar la ilusión de la pequeña escala.
\end{warningbox}

\subsection{Criterio de investigación: saber qué vale la pena hacer}

Yang Songlin dice que el verdadero gran reto no es el problema técnico en sí, sino no saber qué hacer. Su método consiste en leer primero todos los artículos del campo que vale la pena revisar, formar un mapa de problemas y luego juzgar qué mecanismo antiguo tiene valor hoy. Por ejemplo, DataNet ya existía en 2021, pero no tenía garantía de eficiencia en hardware; cuando las necesidades de las tareas y los algoritmos paralelos maduran, puede volver a ser un componente utilizable.

\begin{importantbox}{El criterio en la investigación de algoritmos}
La buena investigación en arquitecturas no consiste en combinar módulos al azar, sino en saber qué herramientas existieron en la historia, por qué no tuvieron éxito, qué restricciones han cambiado hoy y cómo usar algoritmos y hardware nuevos para hacer scale up de las ideas antiguas.
\end{importantbox}

\subsection{La implementación de código abierto también es una contribución de investigación}

Flash Linear Attention es importante no solo por el artículo, sino porque permite que investigadores y la industria prueben con facilidad la atención lineal. Si una idea no tiene una implementación utilizable, a la investigación posterior le cuesta darle seguimiento; si la implementación es fácil de usar, la línea técnica se transmite y se valida con más facilidad.

\begin{knowledgebox}{Tres pasos del artículo a la línea de investigación}
El primer paso es la idea: que tenga sentido en lo matemático o en el mecanismo. El segundo paso es la implementación: si puede escribirse como código utilizable, reproducible y optimizable. El tercer paso es el scaling: si se sostiene con modelos más grandes, tareas más realistas y condiciones más favorables para el hardware.
\end{knowledgebox}

\subsection{Resumen de la sección}

La arqueología de algoritmos ofrece un método de investigación: no hay que perseguir solo la buzzword más reciente, sino leer artículos antiguos, entender los mecanismos y juzgar qué ideas pueden volver a brillar con el hardware y las tareas de hoy.
\section{Afinidad con el hardware: el algoritmo debe poder paralelizarse, usar multiplicación de matrices y escalar}

Después de la arqueología de algoritmos, esta sección aborda el último umbral para que una idea llegue a la práctica: el hardware. Si una idea no puede paralelizarse, no puede escribirse como un kernel eficiente ni aprovechar la capacidad de multiplicación de matrices de las GPU modernas, difícilmente llegará a los modelos frontier.

La última línea técnica es la afinidad con el hardware. Yang Songlin (杨松琳) sostiene que un algoritmo no solo debe tener sentido desde el punto de vista de machine learning, sino que también debe poder paralelizarse, escalar e implementarse de forma eficiente en el hardware. Transformer tuvo éxito en su momento no solo por su gran capacidad expresiva, sino también porque tanto attention como FFN contienen abundantes multiplicaciones de matrices, lo que les da una mayor afinidad con el hardware que RNN/LSTM.

\begin{figure}[H]
\centering
\includegraphics[width=0.96\textwidth]{hardware-friendly-algorithm.png}
\caption{Algoritmos afines al hardware: además de tener sentido matemático, deben poder paralelizarse, usar multiplicación de matrices y escalar. Diagrama conceptual propio, elaborado a partir de la conversación de 01:29:50--01:42:23.}
\end{figure}

\begin{knowledgebox}{Lectura de la figura: la investigación actual en arquitecturas es un diseño conjunto de algoritmos y sistemas}
Que sea matemáticamente sólido es solo el primer paso; también se necesita un algoritmo paralelo, como chunkwise; debe poder aprovechar la multiplicación de matrices; debe tener una implementación como kernel, y, por último, debe escalar al tamaño de los modelos grandes. Si falta cualquiera de estos eslabones, al algoritmo le será difícil llegar a los modelos frontier.
\end{knowledgebox}

\subsection{Chunkwise Algorithm y paralelismo}

Esta subsección concreta la afinidad con el hardware en forma de algoritmos paralelos. La atención lineal suele tener una forma recurrente, y la recurrencia por sí misma no se presta al paralelismo masivo en GPU, por lo que es necesario reescribirla en una forma que pueda paralelizarse por bloques.

Linear Attention suele tener una forma recurrente, pero la recurrencia es intrínsecamente difícil de paralelizar. Para entrenarla de manera eficiente en las GPU modernas se requieren algoritmos paralelos como el chunkwise algorithm, que divide las secuencias largas en bloques: conserva la estructura recurrente y, al mismo tiempo, aprovecha el cómputo en paralelo. Su papel es similar al de FlashAttention respecto a Softmax Attention: no cambia el objetivo del modelo, sino que hace que el algoritmo pueda ejecutarse en el hardware.

\begin{importantbox}{La importancia de los algoritmos paralelos}
Sin un algoritmo paralelo eficiente, incluso una variante de attention muy elegante en teoría no puede llegar al entrenamiento a gran escala. La innovación en arquitecturas y la optimización de sistemas no son etapas sucesivas: juntas determinan si una idea puede llevarse a la práctica.
\end{importantbox}

\subsection{La coordinación con el hardware en DeepSeek}

A continuación se toma DeepSeek como ejemplo de coordinación con el hardware. Su valor no está solo en haber elegido sparse, sino en que la forma de implementar sparse se ajusta a FP8, a la multiplicación de matrices y a una ruta de hardware sin softmax.

En la entrevista se considera que DeepSeek da mucha importancia a la coordinación entre hardware y algoritmos. El indexer de DeepSeek Sparse Attention puede calcular el attention score / logit en FP8, eliminar las costosas operaciones softmax/exp y dejar principalmente multiplicaciones de matrices y la selección Top-K. Esto lo acerca más a los principios del hardware y facilita la coordinación con el equipo de infra.

\begin{figure}[H]
\centering
\includegraphics[width=0.94\textwidth]{deepseek-hardware-codesign.png}
\caption{La coordinación con el hardware en DeepSeek: la selección dispersa, FP8 y la multiplicación de matrices sirven en conjunto a la eficiencia. Diagrama conceptual propio, elaborado a partir de la conversación de 01:29:50--01:42:23.}
\end{figure}

\begin{knowledgebox}{Lectura de la figura: la coordinación con el hardware no consiste solo en escribir kernels}
Del lado del algoritmo, Sparse / Indexer / Top-K reducen el cómputo; del lado del hardware, la eficiencia aumenta con FP8, la multiplicación de matrices y la eliminación de las costosas operaciones softmax/exp. El verdadero co-design considera al mismo tiempo los objetivos del algoritmo y los principios del hardware.
\end{knowledgebox}

\subsection{Resumen de la sección}

La investigación futura en arquitecturas deberá cumplir tres condiciones a la vez: ser sólida desde el punto de vista del aprendizaje automático, ser explicable matemáticamente y, en el hardware, poder paralelizarse y usar multiplicación de matrices. De lo contrario, por elegante que sea un algoritmo, difícilmente llegará a los modelos frontier.
\section{Ruta futura: ¿convergerán Linear y Sparse?}

Esta sección examina en conjunto las rutas anteriores. En el programa se menciona una solución ideal para el futuro: combinar Linear Attention y Sparse Attention. Intuitivamente, Linear/Hybrid se encarga de reducir la KV cache y el costo recurrente de la mayoría de las capas, mientras que Sparse se encarga de seleccionar unos pocos token clave cuando se necesita una recuperación global. En teoría, si la selección sparse es precisa, puede sustituir una parte de las capas de Full Attention.

\subsection{Por qué resulta atractiva una ruta combinada}

Esta subsección explica por qué es probable que en el futuro no gane una sola ruta. Linear, Sparse y Full resuelven cada una un cuello de botella distinto, y los sistemas de ingeniería reales suelen combinarlas.

Linear Attention destaca por ahorrar estado y reducir el costo de la mayoría de los token, pero puede sacrificar algunas interacciones globales complejas; Sparse Attention conserva la capacidad de «mirar las posiciones importantes del historial», pero necesita un indexer suficientemente confiable. La ruta combinada busca que ambas se complementen: la mayor parte del tiempo se usa cheap recurrence y, cuando hace falta, sparse retrieval completa la información de largo alcance.

\begin{knowledgebox}{Tres posibles rutas combinadas}
\begin{tabularx}{\textwidth}{p{0.24\textwidth}p{0.34\textwidth}X}
\toprule
Forma de combinación & Mecanismo & Riesgo principal \\
\midrule
Linear + Full & Mayoría de capas lineales, pocas capas globales & Las capas globales todavía pueden limitar la context window. \\
Sparse + Full & Mayoría de capas dispersas, pocas capas globales & La KV cache sigue siendo grande y el indexer debe ser confiable. \\
Linear + Sparse & Estado lineal más recuperación dispersa & Mayor dificultad para la estabilidad del entrenamiento y mayor complejidad de implementación. \\
\bottomrule
\end{tabularx}
\end{knowledgebox}

\subsection{Qué significa «seleccionar con precisión»}

Lo esencial de Sparse Attention no es mirar menos token, sino mirar los correctos aun mirando menos. En el razonamiento de múltiples saltos, la recuperación en documentos largos, las dependencias de código y las trayectorias de agent, la información clave puede aparecer en posiciones muy lejanas. Si el indexer selecciona mal, se ahorra cómputo, pero el modelo responde mal; si selecciona demasiado, se pierde la ventaja de la dispersión.

\begin{warningbox}{Modos de falla de Sparse}
La falla más peligrosa de la atención dispersa no es la lentitud, sino omitir en silencio token clave. En las tareas de contexto largo, omitir una condición temprana o una restricción lejana puede hacer que todo el razonamiento posterior sea erróneo.
\end{warningbox}

\subsection{Resumen de la sección}

Es posible que el futuro de Attention no consista en elegir una sola entre Linear, Sparse y Full, sino en combinarlas. La verdadera dificultad es lograr que la ruta combinada conserve al mismo tiempo capacidad expresiva, bajo costo, adecuación al hardware y estabilidad del entrenamiento.
\section{Cómo pueden los investigadores jóvenes entrar a la investigación en arquitecturas}

Una vez expuestas la trayectoria técnica de Attention y las restricciones de hardware, esta sección regresa la discusión al método de investigación. La investigación en arquitecturas no consiste solo en leer artículos ni únicamente en saber escribir fórmulas; requiere estar cerca de modelos reales, cómputo real y restricciones de hardware reales. Para un investigador joven, la pregunta más importante es cómo entrar a un entorno donde realmente pueda verificarse el comportamiento a escala.

El consejo para los investigadores jóvenes al final del programa es directo: hacer una pasantía en un lab y seguirle el paso a la frontier, porque trabajar en arquitecturas exige cómputo. El consejo suena sencillo, pero tiene razones muy concretas: la investigación en arquitecturas difícilmente puede evaluarse solo con experimentos a pequeña escala, y una idea que funciona en un modelo pequeño no necesariamente funciona en un modelo grande.

\begin{figure}[H]
\centering
\includegraphics[width=0.96\textwidth]{young-researcher-advice.png}
\caption{El camino para entrar a la investigación en arquitecturas: arqueología de la literatura, cómputo, un frontier lab y conciencia del hardware son todos indispensables. Diagrama conceptual propio, elaborado a partir de la conversación en 01:42:23--01:43:26.}
\end{figure}

\begin{knowledgebox}{Lectura de la figura: la investigación en arquitecturas necesita estar cerca de las restricciones reales}
Leer artículos antiguos para construir un mapa, encontrar problemas útiles, obtener suficiente cómputo para verificar el comportamiento a escala, entrar a un frontier lab para acercarse a un entorno de entrenamiento real y, después, usar implementaciones de código abierto para que la técnica se difunda. Es un camino que va de la comprensión a la verificación y de ahí al impacto.
\end{knowledgebox}

\subsection{Por qué el cómputo es el boleto de entrada}

Muchos problemas de la investigación en arquitecturas solo se manifiestan a partir de cierta escala. La expresividad, la estabilidad del entrenamiento, el KV cache, el decoding, los kernels, el tamaño de lote y el rendimiento del hardware no pueden verificarse por completo con pequeños experimentos en una sola máquina. Sin cómputo, es muy difícil saber si una arquitectura es una línea de desarrollo con futuro o una ilusión de la escala pequeña.

\subsection{Resumen de la sección}

Entrar a la investigación en arquitecturas exige tanto criterio algorítmico como conciencia de ingeniería y de hardware. Un investigador joven no puede limitarse a leer los artículos más recientes ni a hacer experimentos de juguete; lo más importante es acercarse a las restricciones reales de la frontier.
\section{Términos clave: índice de palabras clave de este episodio}

\begin{longtable}{p{0.25\textwidth}p{0.34\textwidth}p{0.32\textwidth}}
\toprule
Término & Explicación en una frase & Papel en este episodio \\
\midrule
Full Attention & Atención global Softmax estándar & Gran capacidad expresiva, pero alto costo en contextos largos. \\
Linear Attention & Variante de atención con complejidad aproximadamente lineal & Línea que exploran a fondo Kimi, Qwen y otros. \\
Sparse Attention & Solo una parte de los token participa en la atención & Línea que explora a fondo DeepSeek. \\
KV Cache & Caché que guarda las key/value históricas durante la inferencia & Cuello de botella de memoria de GPU en contextos largos. \\
KDA & Kimi Delta Attention & Módulo central de atención lineal de Kimi Linear. \\
Delta Rule & Una familia de reglas de actualización recursiva & Antecedente histórico de KDA y Gated DeltaNet. \\
Gated DeltaNet & Variante de DeltaNet con compuertas & Una de las bases de KDA. \\
Decay & Tasa de decaimiento de la memoria & Un decay de grano fino mejora el uso del hidden state. \\
Hybrid Attention & Arquitectura que combina Full con Linear/Sparse & Línea práctica actual. \\
Scaling Ladder & Validar una arquitectura aumentando la escala por etapas & Mecanismo de evaluación interno de Kimi. \\
MoE & Mixture of Experts, mezcla de expertos & Avance arquitectónico ya comprobado del lado de la FFN. \\
FLOPs & Cantidad de operaciones de punto flotante & Sirve para comparar costo de cómputo y eficiencia. \\
NoPE / RoPE & Sin codificación posicional / codificación posicional rotatoria & Relacionado con detalles de entrenamiento y con el modelado de la posición. \\
Chunkwise Algorithm & Algoritmo paralelo por bloques & Permite entrenar en paralelo los algoritmos de tipo recurrence. \\
Hardware-friendly & Adecuado para el hardware & Factor decisivo para que un algoritmo pueda implementarse a gran escala. \\
\bottomrule
\end{longtable}

\subsection{Resumen de la sección}

Este episodio tiene una densidad de términos muy alta, pero la relación central es clara: los contextos largos elevan el costo de Full Attention, y las arquitecturas Linear, Sparse e Hybrid intentan reducirlo. Una línea con verdadero futuro tiene que equilibrar a la vez capacidad expresiva, eficiencia y adecuación al hardware.
\section{Síntesis y ampliación}

\subsection{Conclusiones principales}

\begin{enumerate}
\item En 2025 la innovación en arquitectura vuelve a ser importante porque el crecimiento de los datos se desacelera, aumenta la demanda de contexto largo y el costo de inferencia se vuelve más pesado.
\item Full Attention tiene una gran capacidad expresiva, pero en secuencias largas tanto el cómputo como la KV cache resultan costosos.
\item La idea central de Linear Attention es reescribir la atención cuadrática como una recurrence casi lineal o como una actualización de estado.
\item La KDA de Kimi Linear usa un decay de grano más fino para aprovechar mejor la memoria de un hidden state limitado.
\item Kimi opta por Hybrid Linear, DeepSeek por Sparse y MiniMax M2 regresa a Full, lo que muestra que Attention todavía no ha convergido.
\item La FFN ya se transformó con MoE; Attention es la siguiente zona del Transformer que puede reelaborarse.
\item La arqueología de algoritmos permite volver a validar, con hardware nuevo y a nueva escala, mecanismos que quedaron olvidados en artículos antiguos.
\item La afinidad con el hardware es el requisito mínimo de la investigación en arquitectura: debe poder paralelizarse, expresarse como multiplicaciones de matrices, implementarse en kernels y escalar (scale).
\item Una de las ventajas de DeepSeek es su fuerte conciencia de la coordinación entre algoritmos e infra/hardware.
\item Los investigadores jóvenes que quieran trabajar en arquitectura necesitan estar cerca de un frontier lab y de un entorno con capacidad de cómputo; de lo contrario, les será difícil validar el comportamiento a escala (scale).
\end{enumerate}

\subsection{Preguntas abiertas}

Las preguntas abiertas se dejan para el final porque la ruta de Attention aún no ha convergido. Las decisiones de Kimi, DeepSeek y MiniMax no son definitivas, sino apuestas provisionales de distintas empresas bajo distintas restricciones.

\begin{itemize}
\item ¿Podrá una Hybrid Attention como Kimi Linear seguir acercándose al desempeño de Full Attention a mayor escala?
\item ¿Podrá Sparse Attention sustituir de verdad a las capas globales sin perder la capacidad de recuperación de largo alcance?
\item ¿La proporción hybrid de tres a uno se consolidará como práctica de largo plazo o será una solución de transición?
\item ¿Se prolongará la vigencia de Full Attention gracias a que el hardware siga optimizándose?
\item Después de MoE, ¿un avance en Attention se convertirá en el siguiente consenso de arquitectura?
\item ¿Cómo puede la investigación en arquitectura construir un puente de validación confiable entre la pequeña escala de código abierto y la gran escala de código cerrado?
\end{itemize}

\subsection{Lecturas adicionales}

\begin{itemize}
\item Kimi Linear: An Expressive, Efficient Attention Architecture.
\item Informes técnicos sobre DeepSeek Sparse Attention / Native Sparse Attention.
\item Serie FlashAttention: para entender la coordinación entre algoritmos y hardware.
\item Materiales sobre arquitecturas como MoE, RoPE, NoPE, Mamba, DeltaNet y Gated Linear Attention.
\item Blogs relacionados de Hazy Research: modelado eficiente de secuencias y arquitecturas afines al hardware.
\end{itemize}

\begin{importantbox}{Valoración final}
Lo que más vale la pena conservar de EP119 es una actitud de investigación: no preguntar solo «qué attention está más de moda», sino qué cuello de botella resuelve, qué capacidad expresiva sacrifica, cómo se paraleliza en el hardware, si supera el scaling ladder y si puede rendir de forma estable en modelos frontier reales.
\end{importantbox}

\end{document}
